% =====================================================================
% tesi/capitoli/02-supporto-e-salvataggio.tex
% =====================================================================
% copre: docs/01-fondamenta-salvataggio.md
% ---------------------------------------------------------------------

\chapter{Il supporto fisico del salvataggio, e ciò che ne consegue}
\label{cap:supporto}

Un salvataggio di un gioco Game Boy o Game Boy Advance non è un file: è il contenuto di un chip di memoria alloggiato dentro la cartuccia, e diventa un file soltanto nel momento in cui qualcuno lo copia all'esterno. La distinzione appare pedante e non lo è, perché da essa discende ogni proprietà rilevante di tutto ciò che segue: la ragione per cui un salvataggio si perde, la ragione per cui una scrittura non si annulla, e la ragione per cui questo progetto possiede una regola normativa sul backup anziché una raccomandazione.

\section{Le due tecnologie di memoria, e perché la diagnosi cambia}

Nelle generazioni 1 e 2 il chip è una SRAM\footnote{\term{SRAM}, Static Random Access Memory: memoria volatile che conserva il proprio contenuto soltanto finché è alimentata, e che nelle cartucce Game Boy è tenuta viva da una batteria tampone.} da trentadue kibibyte, cioè una memoria volatile che perde il contenuto quando resta priva di alimentazione. Non lo perde in funzione, perché all'interno della cartuccia è saldata una batteria al litio che mantiene alimentato quel solo chip anche a console spenta. Quella batteria è un componente a vita finita, stimata tipicamente in quindici o vent'anni, e la modalità con cui il suo esaurimento si manifesta è la proprietà più importante da capire: il salvataggio non degrada progressivamente, ma scompare integralmente al primo spegnimento successivo. Ogni cartuccia di generazione 1 e 2 ancora in circolazione ha superato di molto quella vita nominale, il che colloca l'intero corpus di quei salvataggi in una condizione di rischio non teorico.

Nella generazione 3 la situazione cambia, e non in modo uniforme fra i titoli. Rubino, Zaffiro e Smeraldo ospitano ancora una batteria, ma non per il salvataggio: quest'ultimo risiede in una memoria flash da centoventotto kibibyte, che non richiede alimentazione per conservare il proprio contenuto. La batteria alimenta invece il \term{real time clock}, l'orologio interno, dal quale dipendono la maturazione delle bacche, il ciclo delle maree e gli eventi vincolati all'ora del giorno. Quando quella batteria si esaurisce il salvataggio permane e l'orologio si arresta, e il gioco espone il noto messaggio sull'orologio interno non funzionante. Rosso Fuoco e Verde Foglia non implementano l'orologio e di conseguenza non alloggiano alcuna batteria \cite{pokefirered}.

Questa differenza fra le due tecnologie chiarisce una confusione diffusa, e la chiarificazione ha valore diagnostico immediato perché cambia la conclusione a cui si giunge partendo dal medesimo sintomo: su una cartuccia di Smeraldo una batteria esaurita non spiega un salvataggio corrotto, mentre su una cartuccia di Cristallo lo spiega interamente. Chi confonde i due casi cerca la causa nel posto sbagliato, e nel caso di Smeraldo la sostituzione della batteria non produce alcun miglioramento del difetto osservato.

\section{Perché il backup è una norma e non un consiglio}

Il progetto possiede una regola normativa, depositata in \file{.claude/rules/hardware-and-perimeter.md}, che esige un backup in doppia copia su due percorsi distinti, verificato leggibile, prima di ogni scrittura. La ragione della sua forma normativa, cioè del fatto che non ammetta eccezioni negoziabili per ragioni di tempo, è che nessuna delle operazioni di cui questo lavoro si occupa è annullabile.

Quando si scrive sulla SRAM di una cartuccia, il contenuto precedente non viene spostato altrove: viene sovrascritto. Non esiste un cestino, non esiste una cronologia delle revisioni, non esiste una copia di sicurezza prodotta dal gioco, con una sola eccezione parziale che il capitolo~\ref{cap:integrita} descrive e che, per come è congegnata, non costituisce una rete di sicurezza utilizzabile a questo scopo. E il valore di ciò che viene sovrascritto non è ricostruibile in linea di principio, non soltanto in pratica: un salvataggio di vent'anni contiene esemplari catturati in circostanze che non si ripetono, con dati che nessuna procedura sa rigenerare.

Esiste poi una seconda ragione, meno evidente e più insidiosa della prima. Una scrittura può riuscire secondo il software che la esegue e non permanere sul chip, circostanza documentata in particolare sulle cartucce dotate di hardware di salvataggio non originale, e il modo di accorgersene non consiste nell'osservare la schermata del gioco ma nel rileggere i byte e confrontarli con quelli che si intendeva scrivere. Da qui la seconda norma del progetto, quella sulla rilettura verificata, la cui formulazione conviene riportare nei termini in cui è stata scritta: una scrittura che nessuno ha riletto non è una scrittura verificata.

\section{Come un salvataggio diventa un file}

Per trasferire il contenuto della cartuccia su un calcolatore occorre un lettore, cioè un dispositivo che si interpone fra la cartuccia e una porta USB e implementa il protocollo del bus della cartuccia. Il progetto impiega un GBxCart RW pilotato da FlashGBX, e la scheda \file{.claude/context/STACK.md} registra la scelta insieme alle alternative escluse e alle ragioni dell'esclusione.

Esistono almeno tre percorsi alternativi, e la loro conoscenza non è accademica perché ciascuno modifica il perimetro di ciò che risulta possibile. Il primo è un lettore commerciale differente, per esempio il GB Operator, che compie la medesima operazione attraverso un software proprio. Il secondo è una console modificata che legge la cartuccia dall'interno, cioè esattamente ciò che il caso di studio del capitolo~\ref{cap:3ds} realizza per le cartucce Nintendo DS. Il terzo è il più interessante ai fini di questo lavoro: la SRAM può essere letta e scritta senza alcun lettore, inducendo il gioco stesso a eseguire codice attraverso il cavo Link, ed è il meccanismo su cui poggia \file{PkSploit} \cite{pksploit}. Quest'ultima via è trattata nel capitolo~\ref{cap:esecuzione} e non costituisce un espediente marginale, perché è la medesima primitiva su cui si fonda il ponte fra generazioni descritto nel capitolo~\ref{cap:ponte}.

Una volta che il contenuto è divenuto un file, la sua estensione dipende dal software che lo ha prodotto: \file{.sav} e \file{.srm} sono le più diffuse e non implicano alcuna differenza di formato. La dimensione, al contrario, è informativa e va usata come primo controllo: trentadue kibibyte per un salvataggio di generazione 1 o 2, centoventotto kibibyte per uno di generazione 3, e qualunque dimensione differente costituisce il primo indizio che il file non è ciò che sembra. Il capitolo~\ref{cap:integrita} mostra un caso reale in cui proprio questo controllo ha risolto una discordanza fra le fonti.
